\documentclass[11pt,a4paper]{article}
\usepackage[T1]{fontenc}
\usepackage{lmodern}
\usepackage[a4paper,margin=25mm]{geometry}
\usepackage{amsmath,amssymb,amsthm}
\usepackage{xcolor}
\usepackage{microtype}
\usepackage{fancyhdr}
\usepackage[hidelinks]{hyperref}
\hypersetup{pdftitle={Powerful values of 2z cubed plus 1},pdfsubject={Elementary reductions and conditional finiteness}}
\definecolor{ink}{HTML}{18344A}
\newtheorem{proposition}{Proposition}
\newtheorem{theorem}{Theorem}
\newtheorem*{conjecture}{The abc conjecture}
\newcommand{\rad}{\operatorname{rad}}
\newcommand{\Z}{\mathbb Z}
\pagestyle{fancy}
\fancyhf{}
\fancyhead[L]{\small\textcolor{ink}{Powerful values of a cubic}}
\fancyhead[R]{\small $x^3y^2=2z^3+1$}
\fancyfoot[C]{\small\thepage}
\setlength{\headheight}{14pt}
\setlength{\parindent}{0pt}
\setlength{\parskip}{6pt}
\renewcommand{\headrulewidth}{0.3pt}
\begin{document}
\thispagestyle{plain}
\begin{center}
{\LARGE\bfseries\color{ink} Powerful values of $2z^3+1$}\par
\vspace{5pt}
{\large Elementary reductions and conditional finiteness}\par
\vspace{5pt}
{\small The integer equation $x^3y^2=2z^3+1$}
\end{center}
\vspace{5pt}
\textbf{Scope.} An unconditional finiteness proof is not established in this
note. What follows is a self-contained reduction of the question, elementary
necessary conditions, and a complete finiteness proof \emph{assuming the
$abc$ conjecture}. That conditional theorem does not answer the requested
unconditional claim. No internet sources were consulted; no claim of novelty
or of a complete solution list is made.

\section*{1. The exact arithmetic issue}
A positive integer $n$ is called \emph{powerful} if every prime dividing $n$
has exponent at least two in its prime factorization. The integer $1$ is
powerful by this definition. Since $2z^3+1$ is odd, it is never zero for
$z\in\Z$.

\begin{proposition}\label{prop:powerful}
For a fixed integer $z$, there exist integers $x,y$ satisfying
\begin{equation}\label{eq:main}
 x^3y^2=2z^3+1
\end{equation}
if and only if $|2z^3+1|$ is powerful. Each such $z$ gives only finitely
many pairs $(x,y)$.
\end{proposition}
\begin{proof}
Put $n=|2z^3+1|$. If \eqref{eq:main} holds, then $x,y\ne0$, and the exponent
of a prime $p$ in $n$ is $3\alpha+2\beta$, where $\alpha,\beta\ge0$ are its
exponents in $|x|,|y|$. A positive integer of that form is at least two.
Thus $n$ is powerful.

Conversely, write $n=\prod_p p^{e_p}$, where every displayed $e_p\ge2$.
Set $\delta_p=0$ when $e_p$ is even and $\delta_p=1$ when it is odd. In
the odd case $e_p\ge3$, so the integers
\[
 a=\prod_p p^{\delta_p},\qquad
 b=\prod_p p^{(e_p-3\delta_p)/2}
\]
are well defined and satisfy $n=a^3b^2$. Taking
$x=\operatorname{sgn}(2z^3+1)a$ and $y=\pm b$ proves existence.
For $n=1$ the products are empty and equal to $1$.
Finally, any solution has $|x|^3\le n$ and $|y|^2\le n$, so there are
only finitely many pairs.
\end{proof}

The construction makes $|x|$ squarefree. It is the unique representation
with $|x|$ squarefree and $|y|>0$, apart from the sign choice for $y$:
reducing $e_p=3\alpha+2\beta$ modulo $2$ forces
$\alpha=\delta_p$ whenever $\alpha\in\{0,1\}$.

Consequently, the original finiteness question is \emph{exactly} the question
whether $|2z^3+1|$ is powerful for only finitely many integers $z$.
This is a reformulation of the obstacle, not its resolution.

\newpage
\section*{2. Unconditional restrictions and verified solutions}
\begin{proposition}\label{prop:restrictions}
Every integer solution of \eqref{eq:main} satisfies
\[
 x,y\ \text{odd},\qquad
 \gcd(|x|,|z|)=\gcd(|y|,|z|)=1,
\]
and
\[
 z\not\equiv1\pmod3,\qquad y\equiv\pm1\pmod9.
\]
Also, $x>0$ when $z\ge0$, while $x<0$ when $z\le-1$.
\end{proposition}
\begin{proof}
The right side of \eqref{eq:main} is odd, so both $x$ and $y$ are odd.
If a prime divides $z$ and either $x$ or $y$, reducing the equation modulo
that prime gives $0=1$, a contradiction. The sign assertion follows because
$y^2>0$ and $2z^3+1$ is positive exactly when $z\ge0$.

The cubes modulo $9$ are $0,1,-1$, so the possible residues of the right
side are $1,3,-1$. None is $0$. Hence neither $x$ nor $y$ is divisible by
$3$. The possible left-side residues are therefore
\[
 (\pm1)\{1,4,7\}=\{1,2,4,5,7,8\},
\]
which exclude $3$. This rules out $z\equiv1\pmod3$. When $z\equiv0\pmod3$,
the right side is $1$ modulo $9$; this forces $x^3\equiv1$ and
$y^2\equiv1$. When $z\equiv-1\pmod3$, the right side is $-1$; this forces
$x^3\equiv-1$ and again $y^2\equiv1$. Checking the nine residues gives
$y^2\equiv1\pmod9$ exactly for $y\equiv\pm1\pmod9$, as asserted.
\end{proof}

The two values $z=0,-1$ are completely elementary. At $z=0$ the equation is
$x^3y^2=1$, and at $z=-1$ it is $x^3y^2=-1$. Taking absolute values forces
$|x|=|y|=1$. The resulting solutions are exactly
\[
 \boxed{(1,1,0),\ (1,-1,0),\ (-1,1,-1),\ (-1,-1,-1)}
\]
\emph{within those two slices}. No assertion is made that these are all
integer solutions.

\subsection*{An unconditional obstruction to polynomial families}
\begin{proposition}
There are no polynomials $X,Y,Z\in\mathbb Q[T]$ with $Z$ nonconstant
satisfying $X^3Y^2=2Z^3+1$ identically.
\end{proposition}
\begin{proof}
Suppose such an identity exists. Write $a=\deg X$, $b=\deg Y$, and
$c=\deg Z>0$. The polynomials $X,Y$ are nonzero, and comparing degrees
shows $3a+2b=3c$. Also $\gcd(XY,Z)=1$, since any common polynomial divisor
would divide the constant $1$ in the identity.

Differentiation gives the nonzero polynomial
\[
 D=X^2Y(3X'Y+2XY')=6Z^2Z'.
\]
The coprime polynomials $X^2Y$ and $Z^2$ both divide $D$, so their product
divides $D$. Taking degrees yields
\[
 2a+b+2c\le3c-1,\qquad 2a+b\le c-1.
\]
But $c=a+2b/3$, so this would give $a+b/3\le-1$, contradicting
$a,b\ge0$.
\end{proof}

This rules out a polynomial construction with varying $z$; it does not
rule out infinitely many isolated integer solutions. Likewise, finiteness
for each fixed $x$ would leave an infinite union of possible slices.
An unconditional bound on $z$ is still missing.

\newpage
\section*{3. A complete conditional finiteness proof}
For a positive integer $m$, let $\rad(m)$ be the product of its distinct
prime divisors, with $\rad(1)=1$.

\begin{conjecture}
For every $\varepsilon>0$ there is a constant $K_\varepsilon\ge1$ such that
all pairwise coprime positive integers $A,B,C$ with $A+B=C$ satisfy
\[
 C\le K_\varepsilon\,\rad(ABC)^{1+\varepsilon}.
\]
\end{conjecture}

\begin{theorem}[Conditional on $abc$]\label{thm:abc}
Assume the $abc$ conjecture, and put $K=K_{1/11}$. Every integer solution
of \eqref{eq:main} with $z\ne0$ satisfies
\[
 \max\{|2z^3+1|,\,2|z|^3\}\le256K^{11}.
\]
In particular, under this assumption there are only finitely many integer
solutions.
\end{theorem}
\begin{proof}
Let $t=|z|\ge1$, $N=|2z^3+1|$ and $H=\max(N,2t^3)$. For $z>0$ use
\[
 (A,B,C)=(1,2t^3,N),
\]
and for $z<0$ use $(A,B,C)=(1,N,2t^3)$. In both cases $A+B=C=H$.
The two numbers $N$ and $2t^3$ differ by $1$, so the triple is pairwise
coprime, as required.

Proposition~\ref{prop:powerful} implies $\rad(N)\le N^{1/2}$: squaring
the product of its distinct primes gives a divisor of $N$. Also,
$\rad(2t^3)\le2t$. Since $N\le H$ and $t\le(H/2)^{1/3}$, we obtain
\begin{equation}\label{eq:radbound}
 \rad(ABC)\le2t\sqrt N\le2^{2/3}H^{5/6}.
\end{equation}
Applying $abc$ with $\varepsilon=1/11$ now gives
\[
 H\le K\bigl(2^{2/3}H^{5/6}\bigr)^{12/11}
   =K\,2^{8/11}H^{10/11}.
\]
Divide by $H^{10/11}>0$ and raise to the eleventh power to obtain
$H\le2^8K^{11}$. Thus $|z|^3\le128K^{11}$, so only finitely many $z$
are possible. Each admits finitely many pairs $(x,y)$ by
Proposition~\ref{prop:powerful}. The two solutions at $z=0$ were already
handled separately.
\end{proof}

\textbf{Intuition.} Powerful numbers have relatively few distinct prime
factors: their radical is at most their square root. In the consecutive
pair $2t^3,N$, the cubic term contributes at most about $H^{1/3}$ to the
radical and the powerful term at most $H^{1/2}$. Their product is bounded
by a constant times $H^{5/6}$. The exponent $5/6<1$ leaves a gap that
$abc$ converts into a height bound. Without that conjecture, this radical
estimate alone supplies no such bound.

\textbf{Final logical status.} The elementary reductions are unconditional;
the finiteness theorem is conditional. No numerical value of $K$ has
been supplied, and neither an unconditional finiteness proof nor a
complete enumeration is established here.
\end{document}
